\section{Local moments, integral kernels, and the column valuation bound}
\label{cv-section}

The integrality of the final lowering map does not by itself give the
stronger valuation estimate for a column of theta maps.  We prove the
finite algebraic passage from the required theta moments to that
estimate, with the actual determinant factors. The analytic input
is stated explicitly as a condition on the ordinary theta moments;
Section~\ref{tvc-section} assembles its proof for the actual kernel.
In particular, bounded ordinary power moments alone are insufficient.
The source of this part of the construction is the chain of local
kernel estimates and incidence operators in the submitted proof.
\cite{alpo-uniform-source}

\subsection{Polynomial substitution and its determinant factor}

Let $K$ be a number field, let $w$ be a finite prime of $K$, and let
$R=\mathcal O_{K,w}$.  We use the additive valuation $v_w$, with
$v_w(\varpi)=1$ for a uniformizer.  Let $I$ be a finite set.  Fix
$1\le a<b$ and dominant integral weights
$\lambda_\tau\in\mathbb Z^a$ and $\mu_\tau\in\mathbb Z^b$ for
$\tau\in I$.  Write $d_\tau=\lambda_{\tau,a}$ and
$d'_\tau=\mu_{\tau,b}$.  Suppose integers $s_\tau,s'_\tau$ and
nonnegative integers $n_\tau$ satisfy
\begin{align}
 \mu_{\tau,j}+s_\tau&=\lambda_{\tau,j}+s'_\tau
                  &&(1\le j\le a),\label{cv-weight-compatibility}\\
 \mu_{\tau,j}&=-s_\tau&&(a<j\le b),\notag\\
 n_\tau&=d_\tau+s'_\tau.\notag
\end{align}
Thus $d'_\tau=-s_\tau$.  The polynomial partitions underlying these
weights are $\nu_{\tau,j}=\lambda_{\tau,j}-d_\tau$ and
$\nu'_{\tau,j}=\mu_{\tau,j}-d'_\tau$; hence
\[
 \nu'_\tau=(\nu_{\tau,1}+n_\tau,\ldots,
              \nu_{\tau,a}+n_\tau,0,\ldots,0).
 \tag{CV1}
\]
Use the polynomial spaces and integral lattices of
Section~\ref{coef:section}, taking the tensor product over $\tau$.
We write $V_\lambda$ for the polynomial space with column degrees
$\nu_\tau$, and similarly for $V_\mu$.

If $M_\tau=B_\tau R^a$ is a full lattice, define its transported
dual-Weyl lattice $\mathcal N_\lambda(M)$ by the coefficient condition
\begin{equation}\label{cv-lattice-coordinates}
 F\in\mathcal N_\lambda(M)
 \quad\Longleftrightarrow\quad
 F_B(X):=\left(\prod_\tau\det(B_\tau)^{d_\tau}\right)
                       F((B_\tau X_\tau)_\tau)\in R[X].
\end{equation}
This is independent of the bases: a change by an element of
$\operatorname{GL}_a(R)$ and its inverse preserves integral
coefficients, and its determinant is a unit.  The corresponding Weyl
lattice is denoted $\mathcal D_\lambda(M)$; it is contained in
$\mathcal N_\lambda(M)$.  Formula
\eqref{cv-lattice-coordinates} includes the determinant twist even
when $d_\tau$ is negative.

For an $I$-tuple of $a$ by $b$ matrices $A=(A_\tau)$, define
\begin{equation}\label{cv-point-substitution}
 \begin{split}
 F^\sharp(X)&=F(X)\prod_\tau\det(X_\tau)^{n_\tau},\\
 Q_A F(Y)&=F^\sharp\bigl((A_\tau Y_{\tau,[1,a]})_\tau\bigr),
 \end{split}
\end{equation}
where $Y_{\tau,[1,a]}$ consists of the first $a$ columns of the $b$ by
$b$ matrix $Y_\tau$.  The matrices $A_\tau$ may have deficient rank.

\begin{lemma}\label{cv-substitution}
For every $A$, the polynomial $Q_AF$ belongs to $V_\mu$.
If $M'_{\tau}=B'_{\tau}R^b$, $A_\tau M'_\tau\subset M_\tau$,
and $F\in\mathcal N_\lambda(M)$, then
\begin{equation}\label{cv-point-integrality}
 Q_AF\in c(B,B')\mathcal N_\mu(M'),\qquad
 c(B,B')=\prod_\tau\det(B_\tau)^{s'_\tau}
                          \det(B'_\tau)^{-s_\tau}.
\end{equation}
The assertion remains valid after extension to any valued field,
with its valuation ring in place of $R$.
\end{lemma}
\begin{proof}
Multiplication by the determinant powers adds $n_\tau$ to every
column degree of $F$.  Substitution by $A_\tau$ commutes with right
column additions.  The unused columns have degree zero, so all
required column-invariance identities hold, including additions to an
unused column.  Equation (CV1) proves membership in $V_\mu$.

Put $A^0_\tau=B_\tau^{-1}A_\tau B'_\tau$.  The lattice inclusion is
equivalent to $A^0_\tau$ having entries in $R$.  Direct substitution,
using $n_\tau-d_\tau=s'_\tau$, gives
\begin{equation}\label{cv-exact-coordinate-factor}
 \left(\prod_\tau\det(B'_\tau)^{d'_\tau}\right)
                 (Q_AF)((B'_\tau Y_\tau)_\tau)
       =c(B,B')Q_{A^0}F_B(Y).
\end{equation}
The rightmost polynomial has integral coefficients: determinants,
multiplication, and substitution by integral matrices introduce no
denominators.  This proves \eqref{cv-point-integrality} and also proves
the assertion over an arbitrary valuation ring.
\end{proof}

\subsection{Finite moments and passage to a valuation ring}

Replace $A_{\tau,\alpha\beta}$ by independent variables
$Z_{\tau,\alpha\beta}$.  Expand the universal polynomial as
\begin{equation}\label{cv-universal-expansion}
                Q_ZF(Y)=\sum_m Z^m Q_m(F)(Y).
\end{equation}
Every occurring multi-index $m$ satisfies
\begin{equation}\label{cv-finite-degrees}
 \sum_{\alpha=1}^{a}\sum_{\beta=1}^{b}m_{\tau,\alpha\beta}
      =D_\tau:=\sum_{j=1}^{a}(\lambda_{\tau,j}+s'_\tau).
\end{equation}
In particular $D_\tau\ge0$, and the set $\mathcal M$ of such
multi-indices is finite.  For numbers $J_m\in K$, define the linear
operator
\[
                     \mathcal JF=\sum_{m\in\mathcal M}J_mQ_m(F).
 \tag{CV2}
\]
The column-invariance identities and degrees hold coefficient by
coefficient in the independent variables $Z$.  Thus
$\mathcal JF\in V_\mu$ for every choice of the $J_m$.

Fix an embedding $\iota:K\hookrightarrow L$ into a nonarchimedean
valued field such that $|\iota(x)|\le1$ if and only if $x\in R$.
Let $\mathcal O_L$ be its valuation ring.  For a family $\mathscr A$
of tuples of matrices over $L$, suppose each $A\in\mathscr A$
carries $M'_\tau\otimes_R\mathcal O_L$ into
$M_\tau\otimes_R\mathcal O_L$.  For $\kappa\in K$, impose the
following finite approximation condition:
\begin{equation}\label{cv-moment-approximation}
 \begin{gathered}
 \text{for every }\epsilon>0\text{ there are }A_1,\ldots,A_h\in
 \mathscr A\text{ and }c_1,\ldots,c_h\in\mathcal O_L\text{ such that}\\
 \left|\iota(J_m)-\iota(\kappa)\sum_{j=1}^h c_j A_j^m\right|
                  <\epsilon\quad(m\in\mathcal M).
 \end{gathered}
\end{equation}
The same finite sum must approximate every required moment.  Its
length and coefficients may depend on $\epsilon$.

\begin{theorem}\label{cv-moment-integrality}
Under \eqref{cv-moment-approximation},
\[
 \mathcal J\mathcal N_\lambda(M)
       \subset \kappa c(B,B')\mathcal N_\mu(M').
 \tag{CV3}
\]
If $\kappa=0$, then $\mathcal J=0$.
No completeness, compactness, or integral basis for $\mathcal O_L$
is needed in this assertion.
\end{theorem}
\begin{proof}
Fix $F\in\mathcal N_\lambda(M)$.  By
Lemma~\ref{cv-substitution}, each polynomial
\[
 \iota(\kappa)\sum_jc_jQ_{A_j}(\iota F)
\]
belongs to $\iota(\kappa c(B,B'))$ times the target integral
polynomial lattice over $\mathcal O_L$.
After applying the target coordinate transformation in
\eqref{cv-lattice-coordinates}, its coefficients converge to those
of the similarly transformed $\iota(\mathcal JF)$ as
$\epsilon\to0$.  Indeed, each coefficient is a fixed finite linear
combination of the errors in
\eqref{cv-moment-approximation}; the ultrametric inequality bounds
its absolute value by $C_F\epsilon$, where $C_F$ is the maximum of
the finitely many absolute values of its multipliers.

For a nonzero $q\in L$, the ball $q\mathcal O_L$ is closed in the
following elementary sense.  If $x$ can be approximated arbitrarily
closely by $qz$ with $|z|\le1$, choose an approximation with error at
most $|q|$; then $|x|\le|q|$.  For $q=0$, arbitrary approximation
means $x=0$.  Apply this observation coefficient by coefficient with
$q=\iota(\kappa c(B,B'))$.  When $\kappa\ne0$, divide the
coefficients by $q$ and use the defining property of $\iota$ to
descend integrality to $R$.  This proves (CV3).  When $\kappa=0$,
all coefficients vanish, proving the last assertion.
\end{proof}

\begin{lemma}[A bounded measure supplies the finite approximations]
\label{cv-measure-approximation}
Let $X$ be a compact space, let $L$ be complete, and let
$\ell:C(X,L)\to L$ be a continuous linear functional of norm at
most one.  Suppose $a:X\to\mathscr A$ is continuous, and
$J_m=\kappa\ell(a^m)$ for the finitely many relevant $m$.
Then \eqref{cv-moment-approximation} holds, provided $X$ has a
basis of clopen sets.
\end{lemma}
\begin{proof}
For any prescribed error, continuity and compactness give a finite clopen
partition $X=\coprod_{j=1}^h X_j$ such that all the functions $a^m$
oscillate by less than that error on each part.  Such a partition is
obtained by taking a finite clopen subcover and successively removing
earlier sets.  Omit empty parts, choose $x_j\in X_j$, and put
$c_j=\ell(1_{X_j})$.  Since $\|1_{X_j}\|_\infty=1$, one has
$|c_j|\le1$.  Applying $\ell$ to the step-function approximation
gives the required simultaneous estimates.  If $\kappa\ne0$,
divide the prescribed error by $|\kappa|$ first; if $\kappa=0$
the assertion is immediate.
\end{proof}

\subsection{The selected and conjugate primes}

Let $E/F$ be a CM extension with $[E:\mathbb Q]=2f$ and conjugation
$x\mapsto\bar x$.  Assume that a rational prime $p$ and an algebraic
integer $\pi$ satisfy
\[
 \pi\bar\pi=p,\qquad (\pi)+(\bar\pi)=\mathcal O_E,\qquad
 (\pi)=\prod_{P\in\mathcal P}P,\qquad
 |\mathcal P|=f,\quad\mathcal O_E/P=\mathbb F_p.
 \tag{CV4}
\]
Then $p$ is a uniformizer at every $P\in\mathcal P$ and at every
$\bar P$.  Let $N\subset E^a$ and $N'\subset E^b$ be self-dual
lattices for $h(x,y)=\bar x^{\,t}y$.  At every selected prime choose
nonzero linear functionals
\[
 \ell_P:N_P\to\mathbb F_p,\qquad
 \ell'_P:N'_P\to\mathbb F_p,
\]
linear over the local ring through its residue map.  Put
\[
 L_P=\ker\ell_P,\qquad L'_P=\ker\ell'_P,
\]
and let $L,L'$ agree with $N,N'$ at all other primes.  A superscript
$\#$ denotes the Hermitian dual lattice.

\begin{lemma}\label{cv-kernel-lattices}
Suppose $N_P=\mathcal O_{E,P}^a$ and
$N'_P=\mathcal O_{E,P}^b$ for all $P\in\mathcal P$.
Then
\[
 (L^\#)_P=N_P,\quad (L'^\#)_P=N'_P,
 \qquad v_P(\det L_P)=v_P(\det L'_P)=1,
 \tag{CV5}
\]
where determinant valuations are relative to the standard basis.
At conjugate primes,
\[
 v_{\bar P}(\det(L^\#)_{\bar P})
 =v_{\bar P}(\det(L'^\#)_{\bar P})=-1.
 \tag{CV6}
\]
In (CV5) a basis can be chosen with determinant exactly $p$.
\end{lemma}
\begin{proof}
Choose a coordinate $k$ on which $\ell_P$ is nonzero and divide its
residue row by that nonzero coefficient.  Lift the other coefficients
to $a_j\in\mathcal O_{E,P}$.  A basis of its kernel is
\[
                  p e_k,\qquad e_j-a_j e_k\quad(j\ne k).
\]
With the columns in their original order, its determinant is $p$.
The same proof applies to $L'_P$.
Hermitian duality exchanges $P$ and $\bar P$, and in bases sends
$B$ to $\bar B^{-t}$.  Since $L_{\bar P}=N_{\bar P}$ and $N$ is
self-dual, $(L^\#)_P=N_P$.  The inverse determinant gives (CV6).
These arguments use no division by two.
\end{proof}

Define two subsets of $\operatorname{Hom}_E(E^a,E^b)$:
\begin{align}
 \mathscr S_P&=\{x:x^*L'_P\subset N_P
                                \text{ for every }P\in\mathcal P\},
                         \label{cv-supports}\\
 \mathscr S_{\bar P}&=\{x:xN_P\subset N'_P,
                \ \ell'_P(xy)=\ell_P(y)
                \text{ for every }y\in N_P, P\in\mathcal P\}.
                         \notag
\end{align}
Here $x^*=\bar x^{\,t}$.

\begin{lemma}\label{cv-support-duality}
If $x\in\mathscr S_{\bar P}$, then
\[
               x^*(L'^\#)_{\bar P}\subset(L^\#)_{\bar P}
                          \quad(P\in\mathcal P).
 \tag{CV7}
\]
\end{lemma}
\begin{proof}
The functional identity implies $xL_P\subset L'_P$.  For
$z\in(L'^\#)_{\bar P}$ and $y\in L_P$, adjointness gives
$h(y,x^*z)=h(xy,z)\in\mathcal O_{E,\bar P}$.
The local description of the Hermitian dual therefore places
$x^*z$ in $(L^\#)_{\bar P}$.  This description follows directly
in a local basis from $L^\#_{\bar P}=\bar B^{-t}
\mathcal O_{E,\bar P}^a$.
\end{proof}

For the coefficient field $K$, take a set $I$ of $f$ embeddings
$\tau:E\hookrightarrow K$.  At the prime $w\mid p$ suppose either
all primes $\tau^{-1}w$ belong to $\mathcal P$, or all belong to
$\overline{\mathcal P}$.  These are respectively the two cases
below.  Write $e_w=v_w(p)$; for $z\in E$, at either corresponding
prime $Q$ one has
\[
                         v_w(\tau(z))=e_wv_Q(z),
 \tag{CV8}
\]
because $p$ is a uniformizer at $Q$.

\begin{quote}
\textbf{The required ordinary theta moments.}
\label{cv-theta-moments}
For the actual theta transition from rank $a$ to rank $b$, write its
polynomial kernel as $\eta\mathcal J$, where
$\eta\in K^\times$ is its algebraic character factor and
$v_w(\eta)=0$.  The indices used by $\mathcal J$ are exactly
\eqref{cv-finite-degrees}.  There are numbers
$u_\tau\in\mathbb C_p$ with $|u_\tau|=1$ such that its moments
$J_m$ satisfy \eqref{cv-moment-approximation} with matrices
\[
                         A_\tau(x)=u_\tau\,\iota(\tau(x^*))
 \tag{CV9}
\]
and with the following support and scalar:
\[
 \begin{array}{c|c|c}
 \tau^{-1}w\in\mathcal P&x\in\mathscr S_P&\kappa=p^{-fa}\\
 \tau^{-1}w\in\overline{\mathcal P}&x\in\mathscr S_{\bar P}
                                                  &\kappa=1.
 \end{array}
 \tag{CV10}
\]
The approximation is simultaneous for all needed indices, and uses
coefficients of absolute value at most one.
\end{quote}

This hypothesis is a finite consequence of the stronger measure
statement used in the source.  That statement constructs a measure
of norm at most one for the Taylor coefficients of a normalized CM
theta function at prime-to-$p$ torsion translates, after multiplying
the differentiation directions by one common $p$-adic unit period.
It then applies finite translations and the indicated support
conditions.  The measure statement has not been proved merely by
the finite-dimensional arguments above.  It must be supplied with
the actual normalization of the theta function and of its period.

\begin{theorem}[The two local kernel bounds]\label{cv-two-kernels}
Assume $N_P=\mathcal O_{E,P}^a$ and
$N'_P=\mathcal O_{E,P}^b$ for every $P\in\mathcal P$.
Assume the ordinary theta moment condition above and equality of the target
Weyl and dual-Weyl lattices at $w$.  Suppose
$s_\tau-s'_\tau=s_0\ge0$ is independent of $\tau$.
For $F\in\mathcal D_\lambda(L^\#)$, the actual kernel $K=\eta\mathcal J$
satisfies
\begin{align}
 p^{fa+\sum_\tau s_\tau}KF
       &\in\mathcal D_\mu(L')
       &&\text{if }\tau^{-1}w\in\mathcal P,\label{cv-P-bound}\\
 KF&\in p^{fs_0}\mathcal D_\mu(L'^\#)
       &&\text{if }\tau^{-1}w\in\overline{\mathcal P}.
                                                      \label{cv-bar-bound}
\end{align}
\end{theorem}
\begin{proof}
In the selected case the source lattice is standard at every
embedding.  Take its basis to be the identity.  By
Lemma~\ref{cv-kernel-lattices}, choose target bases of $L'$ with
determinant $p$.  Matrices (CV9) carry the target lattice into the
source lattice, by \eqref{cv-supports}; multiplication by $u_\tau$
does not change integrality.  The determinant factor of
Lemma~\ref{cv-substitution} is therefore
\[
                         c(B,B')=p^{-\sum_\tau s_\tau}.
\]
Theorem~\ref{cv-moment-integrality}, with $\kappa=p^{-fa}$,
proves \eqref{cv-P-bound} with target $L'$, first in the dual-Weyl
lattice.  The input Weyl lattice is contained in the source
dual-Weyl lattice, and the assumed target equality finishes the
argument.  The scalar $\eta$ is a unit.

In the conjugate case, Lemma~\ref{cv-support-duality} gives the
required matrix inclusions.  Both source and target dual lattices
have determinant of valuation $-e_w$ in each embedding.  Thus
\[
 v_w(c(B,B'))=e_w\sum_\tau(s_\tau-s'_\tau)=e_wfs_0.
\]
Here $\kappa=1$.  The same theorem gives
$KF\in c(B,B')\mathcal N_\mu(L'^\#)
=p^{fs_0}\mathcal N_\mu(L'^\#)$.  Equality of the two target
lattices gives \eqref{cv-bar-bound}.
\end{proof}

The target equality in this theorem is now supplied by
Corollary~\ref{fh-large-prime-equality}.  Its sufficient condition,
for a target weight $\mu$ of length $b$, is
\[
                       \mu_1-\mu_b+b-1\le p.
 \tag{CV11}
\]
The exact hook-product calculation in
Theorem~\ref{fh-hook-product} proves that
$[H_{\mu-\mu_b},H_{\mu-\mu_b}]$ is then a $p$-adic unit.  The
normalized pairing is integral on the dual-Weyl lattice itself,
which is therefore contained in its pairing dual, the Weyl lattice.
The reverse inclusion always holds.  The actual submitted proof
assumes the stronger inequality $\mu_1-\mu_b+b<p$, so its precise
input is covered.  No equality at arbitrary primes is asserted.

\subsection{Reduction to standard local lattices}

The standard-lattice assumption in Lemma~\ref{cv-kernel-lattices}
does not restrict the global class of a self-dual lattice.
We give the approximation argument needed for that assertion.

\begin{lemma}\label{cv-unitary-approximation}
Let $\mathcal P$ be a finite set of primes of a CM field $E$ such
that no member is the conjugate of another.  For a full lattice
$N\subset E^n$ there exists $g\in\operatorname{GL}_n(E)$ with
$g^*g=1$ and
\[
                     (gN)_P=\mathcal O_{E,P}^{n}
                          \quad(P\in\mathcal P).
 \tag{CV12}
\]
If $N$ is self-dual, the same equality holds at every $\bar P$.
\end{lemma}
\begin{proof}
At a split place of $F$ below $P$, a unitary matrix is determined by
an arbitrary matrix $G_P\in\operatorname{GL}_n(E_P)$; its conjugate
component is $(G_P^*)^{-1}$.  Choose $G_P$ carrying $N_P$ to the
standard lattice.  The property of carrying this lattice to the
standard lattice is open: after composing with $G_P^{-1}$, it is
membership in $\operatorname{GL}_n(\mathcal O_{E_P})$.

We prove simultaneous approximation by rational unitary matrices.
The Cayley transform
\[
              C(A)=(1-A)(1+A)^{-1},\qquad A^*=-A,
 \tag{CV13}
\]
is unitary wherever $1+A$ is invertible.  Its inverse is
$A=(1-G)(1+G)^{-1}$ on $\det(1+G)\ne0$.
At a split place the set of $G$ with $\det(1+G)\ne0$ is dense in
$\operatorname{GL}_n(E_P)$.  To see this, replace $G$ by $zG$ with
$z$ arbitrarily close to $1$; the polynomial $\det(1+zG)$ has only
finitely many roots and is not identically zero.  Thus we may choose
the local $G_P$ in this set while preserving its lattice property.

Approximate the entries of their Cayley inverses simultaneously by
entries of a matrix $B$ over $E$, with $B$ close to the prescribed
$A_P$ at $P$ and close to zero at every $\bar P$.  This is the usual
finite-prime approximation obtained by the Chinese remainder
theorem after clearing denominators.  Then $A=B-B^*$ is
skew-Hermitian and is close to $A_P$ at every selected prime.
The determinant $\det(1+A)$ is nonzero when the approximation is
sufficiently close.  The rational matrix $g=C(A)$ then has the
required local lattice property.  All continuity statements concern
finitely many entries and determinants.  This proof works also in
residue characteristic two, since its fields have characteristic zero.

For the last assertion, local Hermitian duality gives
$(gN)_{\bar P}=((gN)_P)^\#$, with conjugate coefficients.
The standard lattice is its own corresponding dual.
\end{proof}

\begin{corollary}[Transport to arbitrary representatives]
\label{cv-kernel-transport}
Suppose the ordinary theta moment condition holds for all representatives
which are standard at the selected primes.  Suppose also that the actual
kernel is equivariant under rational unitary changes of representatives:
if $g,g'$ act on its source and target lattices, and $T_g,T_{g'}$ are
the induced coefficient-space isomorphisms, then
\[
 T_{g'}K_{N,N'}=K_{gN,g'N'}T_g.
 \tag{CV13a}\label{cv-kernel-equivariance}
\]
Under the other hypotheses of Theorem~\ref{cv-two-kernels}, its two
conclusions hold for arbitrary self-dual $N,N'$.
\end{corollary}
\begin{proof}
Choose $g,g'$ by Lemma~\ref{cv-unitary-approximation}.
The changes of representatives preserve
\eqref{cv-supports}: replace $x$ by $g'xg^{-1}$, the two lattices
by $gN,g'N'$, and the residue functionals by composition with the
inverse.  Adjointness uses $g^{-1}=g^*$, and coefficient lattices
transform by the actions already proved in
Section~\ref{coef:section}.  In particular $T_g$ sends the source
Weyl lattice onto the transformed source Weyl lattice, and $T_{g'}$
sends each target Weyl lattice onto its transformed version.
Apply Theorem~\ref{cv-two-kernels} to $gN,g'N'$ and pull the two
inclusions back through $T_{g'}$, using \eqref{cv-kernel-equivariance}.
\end{proof}

\subsection{Algebraic character factors and stabilizer denominators}

\begin{lemma}\label{cv-character-unit}
Let $\chi$ be a Hecke character with conductor prime to
$p$, with values in a number field $K$, such that a positive power
of $\chi$ is algebraic.  An ideal or finite-idele
argument which is a unit at every prime over $p$ has character value
a unit at every $w\mid p$.  Finite-order character factors always
have this property.  The same assertion holds for a character on
norm-one fractional ideals $I\bar I=1$ whose value on a principal
norm-one ideal $(z)$ is a product of integral powers of the embeddings
of $z$.
\end{lemma}
\begin{proof}
First replace $\chi$ by its stated positive power; vanishing of
the valuation of a power implies vanishing of the original valuation.
For an ideal prime to the conductor and to $p$, raise its ray class
to a positive power which is principal, say generated by $z$ with
$z\equiv1$ at the conductor.  Since the ideal is prime to $p$,
$z$ is a unit at every prime over $p$.  The defining principal-ideal
formula for an algebraic Hecke character expresses this power of
its value as a product of integral powers of the conjugates of $z$.
All those factors are $w$-adic units.  Taking valuations shows that
the original value is also a unit.  Changing the ideal representative
or inserting local unit factors introduces only the same principal
formula and finite-order factors.  A root of unity has valuation
zero because a positive power is one.  This proves the assertion
for finite ideles as well.  The finiteness of the ray class group
used here follows from finiteness of the ordinary class group and
of the residue unit group at the conductor.

For a norm-one ideal $I$, choose $h>0$ with $I^h=(a)$.
Then $I^{2h}=(a/\bar a)$ and
$(a/\bar a)\overline{(a/\bar a)}=1$.  If $I$ is a unit above $p$,
then $a$ and $\bar a$ are units there.  Apply the stated
principal-ideal formula to $I^{2h}$ and take valuations.
\end{proof}

In the actual standard representatives the determinant ideals of
$N,N'$ are units above $p$.  The extra level primes are disjoint
from the selected primes, and the finite auxiliary twists are
unramified there.  The possibly half-integral unitary character types
become algebraic after squaring; equivalently their local values
satisfy the positive-power principal formulas just used.  The
character scalar $\eta$ is a ratio of the
corresponding determinant and local-frame character values.
Lemma~\ref{cv-character-unit} therefore gives $v_w(\eta)=0$.
This verification uses the prime-to-$p$ level condition; it cannot
be inferred solely from the absolute value of a complex character.

\begin{lemma}\label{cv-stabilizer-unit}
Let $\Gamma$ be a finite subgroup of $\operatorname{GL}_n(E)$, where
$[E:\mathbb Q]=d$.  Every prime divisor $p$ of $|\Gamma|$ satisfies
$p\le dn+1$.  Hence the stabilizer orders occurring at ranks at
most $N_m$ are units at $w\mid p$ whenever
\[
                              p>2fN_m+1.
 \tag{CV14}
\]
The same conclusion holds for stabilizers which inject into those
matrix groups.
\end{lemma}
\begin{proof}
Cauchy's theorem gives an element $g$ of order $p$.  Regard it as
a rational linear map on the $dn$-dimensional underlying vector
space.  Its minimal polynomial divides $X^p-1$ and is not $X-1$.
The polynomial $\Phi_p(X)=1+X+\cdots+X^{p-1}$ is irreducible over
$\mathbb Q$: after replacing $X$ by $X+1$, Eisenstein's criterion
at $p$ applies.  Thus the minimal polynomial has degree at least
$p-1$, and $p-1\le dn$.

The stabilizers of the decorated positive definite lattices are
finite by Section~\ref{al-section}.  Their projection to the rational
unitary matrices is injective: once the rational matrix and the
decorated lattice are fixed, every auxiliary level coordinate of a
stabilizing element is determined by its action on the prescribed
frame or residue marking.  This is the stabilizer identification of
Sections~\ref{acl-section} and~\ref{acc-section}.  Apply the first
assertion with $d=2f$ and $n\le N_m$.
\end{proof}

Thus a class-sum kernel with denominators $|\Gamma_c|$ preserves
any common local lattice containment satisfied by its individual
terms.  Multiplication by those denominators costs no valuation at
the selected prime.  This is stronger at that prime than the general
denominator bound in Section~\ref{cp-section}.

\subsection{Incidence sums and removal of one residue condition}

We give the local lattice argument that produces the useful gain.
For a local point $q=(N,\ell)$, let $L(q)=\ker\ell$.  An upper
neighbor $q^+=(N^+,\ell^+)$ has
\[
          N\subset N^+,\qquad L(q^+)=N,
 \tag{CV15}
\]
with the prescribed residue normalization.  The upper-neighbor
sum is an ordinary finite sum over these points.  The degree-one
residue field and the chosen normalization make the neighbor set
finite; the argument below needs no division by its cardinality.

For a subset $V\subset\mathcal P$, let $\mathcal L_i(V,q)$ be the
coefficient lattice obtained by using $L(q)_P$ at $P\in V$ and
$N(q)_P$ at $P\notin V$, with all other local conditions unchanged.
At conjugate primes use the dual lattices.  Fix a coefficient prime
$w$ of the selected type.  The map $\tau\mapsto P(\tau)=\tau^{-1}w$
is a bijection from $I$ to $\mathcal P$.  Indeed, (CV4) makes $p$
split completely in $E$, so the completion at every such prime is
$\mathbb Q_p$.  Two embeddings into $\mathbb C_p$ inducing the same
prime extend to the same embedding of $\mathbb Q_p$, and therefore
coincide.  The embeddings in $I$ are distinct and both sets have
cardinality $f$.

\begin{lemma}\label{cv-incidence-lattice}
At the fixed selected-type coefficient prime $w$, if $P\in V$,
replacing $q$ by a neighbor at $P$ gives
\[
             \mathcal L_i(V,q^+)=\mathcal L_i(V\setminus\{P\},q).
 \tag{CV16}
\]
Consequently the upper-neighbor sum at $P$ sends sections integral
in the first family of lattices to sections integral in the second.
\end{lemma}
\begin{proof}
At $P$ the equality is precisely $L(q^+)_P=N(q)_P$ from (CV15).
Every other selected-prime component is unchanged.  The conjugate
component of the global self-dual neighbor can also change, but none
of the embeddings defining the coefficient lattice at this $w$
sees a conjugate prime.  Transported polynomial lattices localized
at $w$ depend only on the indicated selected-prime components,
giving (CV16).
Each summand therefore lies in the same target lattice, and their
finite sum does also.  The equivariance factors used to express the
sum in class representatives preserve these transported lattices.
\end{proof}

\begin{lemma}\label{cv-erase-cost}
Let $M^0\subset M$ be a local sublattice of residue index $p$,
and let $\lambda$ be a dominant integral weight of length $n$.
Then
\[
                  p^{\lambda_1}\mathcal D_\lambda(M^0)
                              \subset\mathcal D_\lambda(M).
 \tag{CV17}
\]
For tensor factors at $w$, the exponent is the sum of
$\lambda_{\tau,1}$ over the embeddings which see that prime.
The statement permits negative $\lambda_1$.
\end{lemma}
\begin{proof}
Choose a basis in which $M=R^n$ and
$M^0=\operatorname{diag}(p,1,\ldots,1)R^n$; a unit determinant
change is harmless.  The round action on a row-weight component
of polynomial row degree $r_1$ multiplies it by
$p^{-d-r_1}$, where $d=\lambda_n$.
Every monomial of $V_{\lambda-d}$ has row degree
$r_1\le\lambda_1-d$, by the weight-majorization proof of
Section~\ref{am-section}.  Hence
$\lambda_1-d-r_1\ge0$.
The Weyl lattice is the sum of its integral row-weight components:
its generators are iterated divided row operators applied to the
highest vector, each of which has a definite row weight.  Multiplying
each component by the indicated nonnegative power of $p$ preserves
the Weyl lattice.  This proves (CV17).  Multiplying the powers in
the different tensor factors proves the last assertion.
\end{proof}

\subsection{The one-step bound on the actual class lattices}

Consider the transition from rank $N_i$ to $N_{i+1}$, with
$N_{i+1}-N_i=t+1=2s_0$.  Let $\mathscr S_i$ be its local lattice
of equivariant class sections.  The original coefficient lattice at
each point is obtained from its $L^\#$, as above.  Fix a set
$\mathcal T\subset I$ of distinguished embeddings and let
$V_w\subset\mathcal P$ be the corresponding selected primes.
Write $U_i$ for the product of their upper-neighbor sums at rank
$N_i$.  Operators at different primes commute, since they change
independent local components.

Suppose nonnegative integers $W_\tau$ satisfy, for every step,
\begin{equation}\label{cv-weight-bounds}
 0\le s_{i,\tau}\le W_\tau,\qquad
 (\mu_{i,\tau})_1\le W_\tau\quad(\tau\notin\mathcal T),
 \qquad N_i\le N_m.
\end{equation}
Assume the moment hypothesis and target Weyl/dual-Weyl equality
at each step, as well as \eqref{cv-kernel-equivariance} and (CV14).
Then the class-sum operator $K_i$
satisfies
\begin{equation}\label{cv-one-step}
 p^e U_{i+1}K_i\mathscr S_i\subset\mathscr S_{i+1},
 \qquad e=2\sum_\tau W_\tau+fN_m
\end{equation}
at selected-type primes.  At conjugate-type primes it satisfies
$K_i\mathscr S_i\subset\mathscr S_{i+1}$.

\begin{proof}
By \eqref{cv-P-bound}, multiplying the point kernel by
$p^{fN_i+\sum s_{i,\tau}}$ puts its value in the coefficient
lattice having every selected residue condition imposed.  The
class-sum denominator is a unit by
Lemma~\ref{cv-stabilizer-unit}; sums retain this containment.
Apply each upper-neighbor sum at a prime of $V_w$.
Lemma~\ref{cv-incidence-lattice} removes exactly that residue
condition.  For all remaining primes apply
Lemma~\ref{cv-erase-cost}.  The additional exponent is
$\sum_{\tau\notin\mathcal T}(\mu_{i,\tau})_1$; if some terms
are negative, the same inclusion holds and the upper bound below
only becomes easier.  The total exponent is at most
\[
 fN_i+\sum_\tau s_{i,\tau}
       +\sum_{\tau\notin\mathcal T}(\mu_{i,\tau})_1
                    \le fN_m+2\sum_\tau W_\tau=e.
\]
Any remaining nonnegative power of $p$ preserves the target
lattice.  This proves \eqref{cv-one-step}.
In the conjugate case, \eqref{cv-bar-bound} already gives a positive
power $p^{fs_0}$ times the original target lattice, and the same
stabilizer argument applies.
\end{proof}

\subsection{Accumulation along the column}

The preceding estimates must be combined with the correct direction
of the incidence commutation law.  Set $r=|V_w|$ and suppose
\begin{equation}\label{cv-commutation}
                         U_{i+1}K_i=p^{s_0r}K_iU_i.
\end{equation}
This is the finite incidence identity of the theta kernels; it is
an additional identity, not a consequence of their integral moments.
Let
$C_i=K_{i-1}\cdots K_0$, with $C_0=1$.

\begin{proposition}\label{cv-column}
Assume \eqref{cv-one-step} and \eqref{cv-commutation} for
$0\le i<m$.  Then
\begin{equation}\label{cv-column-bound}
 p^{b_i}C_iU_0^i\mathscr S_0\subset\mathscr S_i,\qquad
 b_i=ie+s_0r\frac{i(i+1)}2.
\end{equation}
In particular, if $A\ge0$ and
$\mathscr N=p^AU_0^m\mathscr S_0$, then for any $c\in K$ with
\[
                          v_w(c)\ge e_wb_m
 \tag{CV18}
\]
one has $cC_m\mathscr N\subset\mathscr S_m$.
\end{proposition}
\begin{proof}
Commuting one $U$ through the first $i$ transitions gives
\[
                         U_iC_i=p^{s_0ri}C_iU_0.
 \tag{CV19}
\]
This follows by induction directly from \eqref{cv-commutation}.
The assertion \eqref{cv-column-bound} is trivial for $i=0$.
Apply $p^eU_{i+1}K_i$ to its left side.  By
\eqref{cv-one-step}, the image is integral.  Equations
\eqref{cv-commutation} and (CV19) identify that image as
\[
 p^{b_i+e+s_0r(i+1)}C_{i+1}U_0^{i+1}\mathscr S_0.
\]
Since $b_{i+1}=b_i+e+s_0r(i+1)$, induction proves the formula.
Finally, (CV18) makes $cp^{A-b_m}$ an element of $R$.
Multiply the $i=m$ inclusion by this scalar.
\end{proof}

For the actual set of distinguished embeddings, $r=|\mathcal T|$.
Since $t+1=2s_0$, the final exponent can be written
\begin{equation}\label{cv-final-exponent}
 b_m=mE_p,\qquad
 E_p=2\sum_\tau W_\tau+fN_m
                    +|\mathcal T|\frac{(t+1)(m+1)}4.
\end{equation}
Although the last displayed summand can be a half-integer, $b_m$
is an integer because $m(m+1)$ is even.
The initial lattice in the submitted construction has exactly the
form $p^AU_0^m\mathscr S_0$, with
$A=m\sum_{P\in V_w}b_{\rm top}(\tau(P))\ge0$.
At conjugate-type primes $V_w$ is empty and the needed exponent is
zero.  Thus the finite argument gives the actual column valuation
bound, on its stated initial lattice, once the local moment and
incidence inputs have been supplied.  It does not assert the same
bound on every vector of the original unmodified class lattice.

The final all-twist lowering theorem of
Section~\ref{atl-section} preserves these local integral lattices.
Applying it introduces no further denominator.  Pairing with a row
then uses the finite-orbit integrality theorem of
Section~\ref{cp-section}; its stabilizer denominator is a unit at
$p$ by (CV14).  These two facts connect the local column bound to
the actual period matrix entry, without a stabilizer average in the
lowering map.

\subsection{Counterexamples to weakened assertions and remaining inputs}

The support and normalization conditions above have concrete force.
Here are elementary checks which prevent stronger, false replacements.

\begin{enumerate}
\item Take $a=1$, $b=2$, $M=\mathbb Z_p$, $M'=\mathbb Z_p^2$,
$\lambda=(0)$, $\mu=(1,0)$, $s=0$, and $s'=1$.
These weights satisfy the stated compatibility conditions, with
$n=1$ and the canonical source polynomial $F=1$.
For $A=(p^{-1},0)$, one has
$Q_AF=p^{-1}Y_{11}$, which is not integral.  Thus the matrix
support condition cannot be omitted from the finite moment theorem.
\item Bounding each ordinary power moment by one does not give a
measure of norm at most one on $\mathbb Z_p$.  For example, let
$J_1=1$ and $J_j=0$ for every $j\ne1$, including $j=0$.
The induced functional on polynomials is differentiation at zero.
It takes the integer-valued polynomial $\binom Xp$ to
$(-1)^{p-1}/p$.  Since $|\binom xp|_p\le1$ for every
$x\in\mathbb Z_p$, this functional cannot extend with norm at
most one.  Thus the ordinary theta measure input is materially
stronger than bounded individual derivatives.
\item The coefficient bound $|c_j|\le1$ in
\eqref{cv-moment-approximation} is essential: a single integral
evaluation multiplied by $p^{-1}$ has no such integrality property.
\item For $\lambda=(2,0)$, the polynomial Weyl lattice on one
column is generated by $x^2,2xy,y^2$, whereas the dual-Weyl lattice
contains $xy$.  They differ at $p=2$.  The representation-theoretic
prime bound therefore cannot be discarded.
\item A finite class sum with one class and stabilizer of order $p$
can equal $1/p$, even when its coefficient pairing is integral.
The unit-stabilizer assertion is indispensable for the loss-free
class-sum passage at the selected prime.
\item Reversing \eqref{cv-commutation} changes the column exponent.
For one-dimensional spaces over $\mathbb Q_p$, take $K_i=1$ and
$U_i=p^{s_0ri}$.  Then \eqref{cv-commutation} holds exactly, while
the reversed factor $p^{-s_0r}$ does not if $s_0r>0$.
\end{enumerate}

The finite approximation, substitution, lattice, determinant,
stabilizer, and accumulation arguments in this section are proved
in full. For the submitted theta construction, the normalized ordinary
theta moment theorem and its exact finite translation/support identities
are proved in Sections~\ref{gm-section} and \ref{csp-section}.
Section~\ref{hec-section} proves \eqref{cv-commutation};
Section~\ref{tvc-section} identifies the actual kernel with (CV2)
and assembles the selected-prime column bound.
Weyl/dual-Weyl equality at the target primes satisfying (CV11) has
been proved by the hook-norm argument of
Section~\ref{fh-section}. These inputs are proved independently
of the final branching map.


